\section{The system under study}
\label{sec:system}

Every measurement in this paper is taken through one shipped system, CrossAudit. Every auditor
prompt is built by the product's own prompt builder, and every family but one is read through
the product's own audit call; the exception, \texttt{gpt-6-astra}, receives the same prompt
bytes through a direct call. We describe it first because
its design fixes what the later sections can and cannot vary. We state its structure only here
and quote no measurement of it. Its limits are the subject of
Section~\ref{sec:ceiling} onward.

\textbf{Separation of roles.} A generator model writes an increment (a set of files) into a Git
repository. An auditor model then reads the committed bytes and returns findings, each graded
BLOCKER or ADVISORY. The auditor must come from a different vendor than the generator: the audit
refuses to run when the configuration names the same vendor for both roles, and when the
generator's vendor is not declared, the receipt records that independence could not be asserted.
The same-vendor arms below run that second way, on purpose. The premise is that a model
reviewing its own vendor's output is an interested reader. The product takes that premise as given, and Section~\ref{sec:ceiling} puts
one consequence of it to the test: whether a same-vendor auditor catches more or less.

\textbf{Two tiers of evidence.} Before the auditor reads anything, a deterministic check layer
runs registered checks over the same committed bytes. A check's finding is \emph{verified}: code
produced it and code can reproduce it. A model's finding is \emph{raised}: it is a reading, not
yet reproduced. The checks are grouped into named profiles, so strictness is a setting the
project chooses:
\begin{itemize}\itemsep0pt
\item \texttt{general}, the default: structured files parse, every local path a file lists as an
  input or dependency exists, and relative links and leftover placeholders are reported as
  advisory.
\item \texttt{science}: results carry a schema; quantities carry units; a result that reports
  convergence must meet its own threshold; every input a result depends on is listed, exists and is cited at an exact revision
  (\emph{provenance}); and every number a report states is annotated with the file and quoted span
  it came from, which a check matches against the cited bytes (\emph{number$\to$source}).
\item \texttt{research}: the general pack plus a citation check under which a report may cite only
  sources it fetched through a governed tool.
\end{itemize}
The science profile follows one rule throughout, and the rule is the one Act 4 tests: \emph{the
model names its evidence, and code verifies it.} A number with no source line, a unit that
disagrees with its schema or an input that does not exist fails without any model's opinion.

\textbf{Verdicts are computed, not voted.} The verdict is derived by code from the evidence, in a
fixed order. A failed deterministic check blocks, whatever the model says. A model BLOCKER with
no deterministic support either blocks or escalates to a person, according to a project setting.
A blocked increment goes back to the generator for a bounded number of revision rounds, three by
default. Each revision passes a repair guard before it is committed: an edit outside the audited
scope is refused, and a new catch-all handler, a deleted assertion or a skipped test is flagged to
the auditor for the next round.

\textbf{Receipts.} A passing round produces a receipt that binds the verdict to the commit that was
audited, the commit of the rulebook the auditor was given, and a digest over every finding with its
tier. Every model call is written to a ledger with its prompt hash and cost. A verification command
recomputes the digest from the recorded evidence, so a receipt can be checked by someone who does
not trust the operator.

\textbf{What the rest of the paper measures.} Acts 2 and 3
(Section~\ref{sec:ceiling}, with Appendices~\ref{sec:method} and~\ref{sec:results}) measure the model tier on its own. There,
an increment counts as flagged when the auditor grades any finding BLOCKER, and the general check
pack runs as in the product but decides no recall or false-positive outcome. The model tier is the
part of the system whose limit is least known, since a deterministic check catches exactly what it
was written to catch. Section~\ref{sec:ai4s} measures the model tier on scientific code, and sets
it beside deterministic verification on reported results and on data, asking what each flags that the
other does not.
